\section{System Architecture}
\label{sec:architecture}

\subsection{Design Assumption}
\label{subsec:arch-assumption}

The subsystem is designed under the assumption that the host application does not, and should not need to, know anything about FIDO2's internals. The host calls \texttt{fido2\_init()}, \texttt{fido2\_start()}, and optionally registers a state callback; everything else, transport framing, PIN protocol state, credential storage, attestation signing, happens inside the subsystem's own thread and work queue context. This is the ``side feature on someone else's firmware'' idea we are presenting: the subsystem must coexist with, not restructure, whatever the host application was already doing.

\subsection{Four-Tier Layer Structure}
\label{subsec:arch-fourtier}

The subsystem's modules do not fall into a single fixed/pluggable split. They fall into four tiers, cutting across three independent dimensions: instance arity (single vs.\ multi), interface width (narrow vs.\ wide), and whether the concern is backend-selectable at all. Table~\ref{tab:fourtier} gives the mechanism realizing each tier.

\begin{table*}[t]
\centering
\caption{Four-tier layer structure, by arity, interface width, and mechanism.}
\label{tab:fourtier}
\small
\begin{tabular}{@{}p{0.13\linewidth} p{0.11\linewidth} p{0.05\linewidth} p{0.13\linewidth} p{0.50\linewidth}@{}}
\toprule
\textbf{Tier} & \textbf{Concern(s)} & \textbf{Arity} & \textbf{Interface width} & \textbf{Mechanism} \\
\midrule
Public, narrow, single-instance & Attestation, User Presence & Single & Narrow (1--2 ops) & Bare symbol, no vtable: \texttt{fido2\_attestation\_sign()}, \texttt{fido2\_up\_wait()}/\texttt{fido2\_up\_cancel()} declared directly; backend chosen by which \texttt{.c} file provides the symbol. \\
\addlinespace
Public, wide, single-instance & Storage & Single & Wide ($\sim$15 ops) & \texttt{extern const struct fido2\_storage\_api fido2\_storage\_backend}, linker-resolved, plus a public wrapper layer (\texttt{fido2\_storage\_store()}, \texttt{fido2\_storage\_load()}, \ldots) dispatching into it. \\
\addlinespace
Public, wide, multi-instance & Transport & Multi & Wide (4 ops, multi-inst.) & \texttt{struct fido2\_transport\_api} vtable plus \texttt{FIDO2\_TRANSPORT\_DEFINE()} / \texttt{STRUCT\_SECTION\_ITERABLE}; USB HID and BLE are both real, concurrently-registrable implementations against this API. \\
\addlinespace
Fixed, never backend-selected & Crypto, CBOR, ClientPIN & n/a & n/a & No Kconfig \texttt{choice}, no backend struct; \texttt{PSA\_CRYPTO}/\texttt{ZCBOR}/\texttt{ZCBOR\_CANONICAL} are \texttt{select}ed unconditionally at the top-level \texttt{menuconfig FIDO2}. \\
\bottomrule
\end{tabular}
\end{table*}

\textbf{Falsifiability test (Kconfig).} Only the four backend-selectable concerns, Attestation, Transport, User Presence, and Storage, have subdirectory Kconfig files \texttt{rsource}d from the top-level \texttt{menuconfig FIDO2}: \texttt{attestation/Kconfig}, \texttt{transport/Kconfig}, \texttt{up/Kconfig}, and \texttt{storage/Kconfig}. Crypto and CBOR appear nowhere as a \texttt{choice}; \texttt{PSA\_CRYPTO}, \texttt{ZCBOR}, and \texttt{ZCBOR\_CANONICAL} are \texttt{select}ed unconditionally at \texttt{menuconfig FIDO2} alongside a long list of concrete PSA algorithm/key-type selects (\texttt{PSA\_WANT\_ALG\_ECDSA}, \texttt{PSA\_WANT\_ALG\_SHA\_256}, \texttt{PSA\_WANT\_ALG\_GCM}, \texttt{PSA\_WANT\_ALG\_HMAC}, \texttt{PSA\_WANT\_ALG\_CBC\_NO\_PADDING}, and others), none of which are gated behind a \texttt{choice}.

The Kconfig structure also independently confirms the arity split in Table~\ref{tab:fourtier}. Attestation, User Presence, and Storage are each declared inside a \texttt{choice} block (\texttt{FIDO2\_ATTESTATION\_BACKEND}, \texttt{FIDO2\_UP\_BACKEND}, \texttt{FIDO2\_STORAGE\_BACKEND}), which Kconfig enforces as mutually exclusive, exactly the single-instance property those tiers require. Transport, by contrast, has no \texttt{choice} block at all: \texttt{FIDO2\_TRANSPORT\_USB\_HID} and \texttt{FIDO2\_TRANSPORT\_BLE} are independent \texttt{bool} options and can both be enabled simultaneously, matching the multi-instance, iterable-section mechanism in Table~\ref{tab:fourtier}, since a device may legitimately expose more than one transport at once. This is a second, independent falsifiable signal for the arity column beyond the vtable-versus-bare-symbol mechanism already cited: Kconfig's own exclusivity semantics, not just the C-level registration mechanism, track the same single/multi split.

%TODO before submission: FIDO2_STORAGE_SECURE in storage/Kconfig and the PSA Protected Storage `depends on` symbol is correct for the targeted Zephyr version.
Storage's \texttt{choice} block, \texttt{FIDO2\_STORAGE\_BACKEND}, currently offers three backends (\texttt{FIDO2\_STORAGE\_SETTINGS}, \texttt{FIDO2\_STORAGE\_NONE}, \texttt{FIDO2\_STORAGE\_CUSTOM}), with a fourth, \texttt{FIDO2\_STORAGE\_SECURE}, backed by the PSA Protected Storage API, added ahead of the evaluation in Section~\ref{sec:evaluation} to support the SETTINGS-versus-secure-storage cost comparison in RQ3.

One further falsifiable detail worth mentioning: \texttt{FIDO2\_UP\_ALWAYS}, the ``always approve'' user-presence backend, \texttt{select}s \texttt{NOT\_SECURE} in the Kconfig tree itself. The security warning of bypassing user presence is not only documented in prose, it is also encoded as a build-time-visible dependency a reviewer of the Kconfig tree would see without reading any documentation.

\subsection{Crypto as a Fixed, Delegated Tier}
\label{subsec:arch-crypto}

\texttt{fido2\_crypto.c} is a fixed wrapper, not a backend selected via a struct, and has no Kconfig choice at all. Portability is delegated one level down to the PSA Crypto API~\cite{psa-crypto-api}, which abstracts hardware-backed versus software crypto beneath the subsystem. This is portability-via-delegation, distinct from and weaker than the portability-via-backend-selection that Storage, Attestation, User Presence, and Transport get through Kconfig. \texttt{fido2\_crypto.h} and \texttt{fido2\_cbor.h} both use plain local include guards (\texttt{FIDO2\_CRYPTO\_H\_}, \texttt{FIDO2\_CBOR\_H\_}) with no \texttt{@ingroup} or version tag, unlike the four public headers described below, further confirming that these are internal implementation files split for maintainability rather than pluggability.

\subsection{Minimal Public Surface}
\label{subsec:arch-surface}

The entire application-facing API is five functions and one callback type: \texttt{fido2\_init()}, \texttt{fido2\_start()}, \texttt{fido2\_stop()}, \texttt{fido2\_set\_state\_callback()}, \texttt{fido2\_get\_state()}, and \texttt{fido2\_reset()}. Nothing else crosses the host/subsystem boundary; a host application that only ever calls these six symbols has no way to observe or alter subsystem-internal state.

\subsection{Mutual Ignorance Between Modules}
\label{subsec:arch-ignorance}

This is the concrete property the composability claim rests on, and it is falsifiable rather than qualitative: grepping for \texttt{\#include "fido2\_crypto.h"} outside \texttt{fido2\_core.c}, \texttt{fido2\_clientpin.c}, and the crypto module itself should return zero hits, and the same check generalizes across every module pair below.

\begin{itemize}
\item \textit{Transport doesn't know about PSA crypto.} The USB HID and BLE transport implementations operate purely on raw bytes, channel identifiers, and their respective wire framing (CTAPHID packet reassembly for USB HID); neither includes a crypto header.
\item \textit{Crypto doesn't know about transport.} \texttt{fido2\_crypto.c} operates purely on buffers and PSA key handles.
\item \textit{Neither transport nor crypto knows about CBOR.} \texttt{fido2\_cbor.c} is the only module that understands CBOR structure; \texttt{fido2\_storage.h} and \texttt{fido2\_attestation.h} include only \texttt{fido2\_types.h}, never \texttt{fido2\_cbor.h}.
\item \textit{Storage doesn't know about crypto or attestation.} Storage backends persist opaque \texttt{struct fido2\_credential} blobs with no dependency on PSA key types or attestation formats.
\item \textit{Attestation doesn't know about transport or storage.} \texttt{fido2\_attestation\_sign()} consumes only \texttt{auth\_data}, \texttt{client\_data\_hash}, and a key identifier.
\item \textit{User presence doesn't know about anything else.} \texttt{fido2\_up\_wait()}/\texttt{fido2\_up\_cancel()} is a pure blocking-gesture abstraction with a single kernel-header dependency.
\end{itemize}

Only \texttt{fido2\_core.c} is permitted to include all of the above; it is the single point of coupling in the subsystem.

\subsection{Zero-Touch Extensibility}
\label{subsec:arch-extensibility}

Adding a new transport requires one new source file implementing the four-function \texttt{fido2\_transport\_api} (\texttt{init}, \texttt{send}, \texttt{notify}, \texttt{shutdown}), one call to \texttt{FIDO2\_TRANSPORT\_DEFINE()}, and one Kconfig option, with nothing in \texttt{fido2\_core.c} changing. This is not a hypothetical: a Bluetooth Low Energy transport was subsequently authored against this exact interface, registering via the same \texttt{STRUCT\_SECTION\_ITERABLE} mechanism as the USB HID transport, with zero modifications to \texttt{fido2\_core.c}, \texttt{fido2\_cbor.c}, or \texttt{fido2\_clientpin.c}. Section~\ref{sec:evaluation} reports the exact diff.

\subsection{Static Allocation, No Heap}
\label{subsec:arch-static}

Large per-command structures, \texttt{mc\_params}, \texttt{ga\_params}, \texttt{cp\_params}, and the shared transmit buffer \texttt{ctap\_tx\_frame[]}, are file-scope statics in \texttt{fido2\_core.c} rather than stack-allocated, keeping the processing thread's stack (\texttt{CONFIG\_FIDO2\_THREAD\_STACK\_SIZE}) small and its worst-case usage analyzable. The same pattern extends to the receive path: \texttt{rx\_enqueue\_msg} and \texttt{rx\_dequeue\_msg} are separate file-scope buffers, guarded by \texttt{rx\_enqueue\_mutex}, used to move a message into the two-entry \texttt{fido2\_msgq} without a stack-resident copy. No dynamic allocation occurs anywhere in the command-processing path.

\subsection{Scope: Implemented Command Table}
\label{subsec:arch-scope}

Table~\ref{tab:scope} states plainly what is implemented, stubbed, or absent as of this submission, derived directly from the \texttt{switch} statements in \texttt{process\_command()} and \texttt{handle\_client\_pin()} rather than from a design description.

\begin{table}
\centering
\caption{CTAP2.1 command and subcommand implementation status.}
\label{tab:scope}
\small
\begin{tabular}{@{}p{0.52\linewidth} p{0.38\linewidth}@{}}
\toprule
\textbf{Command / Subcommand} & \textbf{Status} \\
\midrule
authenticatorMakeCredential & Implemented \\
authenticatorGetAssertion & Implemented \\
authenticatorGetNextAssertion & Implemented \\
authenticatorGetInfo & Implemented \\
authenticatorClientPIN: getPINRetries & Implemented \\
authenticatorClientPIN: getKeyAgreement & Implemented \\
authenticatorClientPIN: setPIN & Implemented \\
authenticatorClientPIN: changePIN & Implemented \\
authenticatorClientPIN: getPinToken (legacy) & Implemented \\
authenticatorClientPIN: getPinUvAuthTokenUsingPinWithPermissions & Implemented \\
authenticatorClientPIN: getPinUvAuthTokenUsingUvWithPermissions & Stubbed (returns \texttt{INVALID\_SUBCOMMAND}) \\
authenticatorClientPIN: getUVRetries & Stubbed (returns \texttt{INVALID\_SUBCOMMAND}) \\
authenticatorSelection & Implemented \\
authenticatorReset & Stubbed (no-op, returns success) \\
authenticatorCredentialManagement & Stubbed (returns \texttt{INVALID\_COMMAND}) \\
PIN/UV Auth Protocol Two & Unimplemented (Protocol One only) \\
authenticatorLargeBlobs, authenticatorBioEnrollment, authenticatorConfig & Not present \\
\bottomrule
\end{tabular}
\end{table}

Two scope notes worth stating explicitly rather than leaving for a reviewer to discover. First, \texttt{authenticatorReset} currently returns success without performing a reset; this is a functional gap, not a security-relevant one, since it does not misrepresent authenticator state, but it means the command cannot yet be relied on. Second, PIN/UV Auth Protocol Two is entirely unimplemented, not merely undocumented, all three of \texttt{decapsulate()}, \texttt{decrypt()}, and \texttt{encrypt()} reject it explicitly. \texttt{authenticatorGetInfo} correctly advertises only Protocol One, so this is an internally consistent scope limitation rather than a spec violation.

\begin{figure}[t]
\centering
\includegraphics[width=\linewidth]{fig/boundary.pdf}
\caption{Module boundary diagram. Every edge terminates at \texttt{fido2\_core.c}; no direct edges exist between peer modules. Node color indicates the four-tier classification from Table~\ref{tab:fourtier}. Companion to Table~\ref{tab:fourtier}: the table gives the mechanism per tier, this figure gives the topology.}
\label{fig:boundary}
\end{figure}

\begin{figure}[t]
\centering
\includegraphics[width=\linewidth]{fig/runtime-state.pdf}
\caption{Runtime state transition diagram, derived from every \texttt{set\_runtime\_state()} call site in the subsystem. \texttt{IDLE} is the hub state, exposed externally via the ISR-safe \texttt{fido2\_get\_state()}; three of four transitions terminate there. This is the actual runtime state machine, not command dispatch, which is a flat switch statement (Table~\ref{tab:scope}), not a state machine.}
\label{fig:runtime-state}
\end{figure}